\par\Needspace{4\baselineskip}
\begin{center}
\SourceLeaf{031}\pdfbookmark[2]{13.}{article-013}13.
\end{center}
\emph{Beispiel.} Die Exzentrizität sei dieselbe wie im Beispiel in Art. 10, und es sei
\[
M = 332^0\,28^{\prime}\,54^{\prime\prime}{,}76.
\]
Drückt man die Exzentrizität in Bogensekunden aus, so beträgt hier $\log e$ $4{,}704\,1513$, und somit ist $e=50600^{\prime\prime}=14^0\,3^{\prime}\,20^{\prime\prime}$. Da $E$ hier kleiner als $M$ sein muss, setzen wir für die erste Rechnung $\epsilon=326^0$ und erhalten aus den kleineren Tafeln
\begin{center}
\small
\begin{tabular}{@{}llp{13em}@{}}
\makebox[7em][l]{$\log\sin\epsilon$\hfill .\,.\,.\,.\,.} & $9{,}74756_n,$ & Änderung je $1^{\prime}$ .\,.\,.\,.\,. $19$, also $\lambda=0{,}32$\\[.3em]
\makebox[7em][l]{$\log e$ in Sek.\hfill .\,.\,.} & $4{,}70415$ & \\
\cline{1-2}
 & $4{,}45171_n$ &
\end{tabular}
\end{center}
\begin{center}
\small
\begin{tabular}{@{}llp{12em}@{}}
\multicolumn{2}{@{}l}{daraus $e\sin\epsilon=-28295^{\prime\prime}=-7^0\,51^{\prime}\,35^{\prime\prime}$.} & Änderung des Logarithmus je Tafeleinheit, die hier $10$\\
$M+e\sin\epsilon=$ & $324\ \ 37\ \ 20$ & Bogensekunden entspricht, .\,.\,.\,.\,. $16$, also $\mu=1{,}6$\\[.3em]
Abweichung von $\epsilon$ .\,.\,.\,.\,.\,.\,.\,.\,.\,.\,.\,. & $1\ \ 22\ \ 40=4960^{\prime\prime}$. &
\end{tabular}
\end{center}
\noindent Daraus folgt $\frac{0{,}32}{1{,}28}\times4960^{\prime\prime}=1240^{\prime\prime}=20^{\prime}\,40^{\prime\prime}$. Der korrigierte Wert lautet somit $E$~$=324^0\,37^{\prime}\,20^{\prime\prime}-20^{\prime}\,40^{\prime\prime}=324^0\,16^{\prime}\,40^{\prime\prime}$. Mit diesem Wert wiederholen wir die Rechnung anhand der größeren Tafeln.
\begin{center}
\begin{tabular}{@{}lll@{}}
\makebox[9em][l]{$\log\sin\epsilon$\hfill .\,.\,.\,.\,.\,.} & $9{,}766\,3058_n$ & $\lambda=29{,}25$\\[.3em]
\makebox[9em][l]{$\log e$\hfill .\,.\,.\,.\,.\,.\,.\,.} & $4{,}704\,1513$ & \\
\cline{1-2}
 & $4{,}470\,4571_n$ & $\mu=147$
\end{tabular}
\end{center}
\begin{center}
\begin{tabular}{@{}ll@{}}
\multicolumn{2}{l}{$e\sin\epsilon=-29543^{\prime\prime}{,}17=-8^0\,12^{\prime}\,23^{\prime\prime}{,}17$}\\[.3em]
$M+e\sin\epsilon=$ & $324\ \ 16\ \ 31{,}59$\\[.3em]
Abweichung von $\epsilon$ .\,.\,.\,.\,.\,.\,.\,.\,.\,.\,.\,.\,.\,.\,.\,. & $8{,}41$.
\end{tabular}
\end{center}
\noindent Multipliziert man diese Differenz mit $\frac{\lambda}{\mu-\lambda}=\frac{29{,}25}{117{,}75}$, so erhält man $2^{\prime\prime}{,}09$. Der erneut korrigierte Wert $E=324^0\,16^{\prime}\,31^{\prime\prime}{,}59-2^{\prime\prime}{,}09=324^0\,16^{\prime}\,29^{\prime\prime}{,}50$ ist damit auf $0^{\prime\prime}{,}01$ genau.
